\documentclass[12pt, a4paper]{report}

% ── Encodage & langue ──────────────────────────────────────────────────────────
\usepackage[utf8]{inputenc}
\usepackage[T1]{fontenc}
\usepackage[french]{babel}

% ── Mise en page ───────────────────────────────────────────────────────────────
\usepackage[top=2.5cm, bottom=2.5cm, left=3cm, right=2.5cm]{geometry}
\usepackage{setspace}
\onehalfspacing

% ── Typographie ────────────────────────────────────────────────────────────────
\usepackage{lmodern}
\usepackage{microtype}

% ── Couleurs & graphiques ──────────────────────────────────────────────────────
\usepackage[dvipsnames, table]{xcolor}
\usepackage{graphicx}
\usepackage{tikz}
\usetikzlibrary{shapes.geometric, arrows.meta, positioning, shadows}

% ── Tableaux ───────────────────────────────────────────────────────────────────
\usepackage{booktabs}
\usepackage{tabularx}
\usepackage{longtable}
\usepackage{multirow}
\usepackage{array}

% ── Code source ────────────────────────────────────────────────────────────────
\usepackage{listings}

% ── Liens & PDF ────────────────────────────────────────────────────────────────
\usepackage{hyperref}
\hypersetup{
  colorlinks  = true,
  linkcolor   = TropicalTeal,
  urlcolor    = TropicalTeal,
  citecolor   = TropicalTeal,
  pdftitle    = {Rapport de Projet — Tropical Medical},
  pdfauthor   = {Yakhoub Baby},
}

% ── Entêtes & pieds de page ────────────────────────────────────────────────────
\usepackage{fancyhdr}
\pagestyle{fancy}
\fancyhf{}
\fancyhead[L]{\small\textit{Tropical Medical — Plateforme Médicale Sénégalaise}}
\fancyhead[R]{\small\textit{Yakhoub Baby}}
\fancyfoot[C]{\thepage}
\renewcommand{\headrulewidth}{0.4pt}

% ── Titres de chapitres ────────────────────────────────────────────────────────
\usepackage{titlesec}
\titleformat{\chapter}[block]
  {\normalfont\huge\bfseries\color{TropicalTeal}}
  {\thechapter.}{12pt}{}
\titleformat{\section}
  {\normalfont\Large\bfseries\color{TropicalDark}}
  {\thesection}{8pt}{}
\titleformat{\subsection}
  {\normalfont\large\bfseries\color{TropicalDark}}
  {\thesubsection}{6pt}{}

% ── Boîtes colorées ────────────────────────────────────────────────────────────
\usepackage{tcolorbox}
\tcbuselibrary{skins, breakable, listingsutf8}

% ── Listes ─────────────────────────────────────────────────────────────────────
\usepackage{enumitem}

% ── Couleurs du projet ─────────────────────────────────────────────────────────
\definecolor{TropicalTeal}{RGB}{14, 159, 142}
\definecolor{TropicalDark}{RGB}{15, 23, 42}
\definecolor{TropicalLight}{RGB}{240, 253, 250}
\definecolor{CodeBg}{RGB}{248, 250, 252}
\definecolor{AccentGold}{RGB}{230, 184, 0}

% ── Style boîte technologie ────────────────────────────────────────────────────
\newtcolorbox{techbox}[2][]{
  colback    = TropicalLight,
  colframe   = TropicalTeal,
  fonttitle  = \bfseries\color{white},
  coltitle   = white,
  colbacktitle = TropicalTeal,
  title      = #2,
  sharp corners = south,
  #1
}

% ── Style boîte remarque ───────────────────────────────────────────────────────
\newtcolorbox{remarque}{
  colback  = yellow!8,
  colframe = AccentGold,
  fontupper = \small,
  left = 4pt, right = 4pt,
}

% ── Configuration listings ─────────────────────────────────────────────────────
\lstset{
  backgroundcolor = \color{CodeBg},
  basicstyle      = \ttfamily\small,
  breaklines      = true,
  frame           = single,
  rulecolor       = \color{TropicalTeal!40},
  keywordstyle    = \bfseries\color{TropicalTeal},
  commentstyle    = \itshape\color{gray},
  stringstyle     = \color{AccentGold!80!black},
  numbers         = left,
  numberstyle     = \tiny\color{gray},
  xleftmargin     = 1em,
  extendedchars   = true,
  literate        =
    {é}{{\'{e}}}1 {è}{{\`{e}}}1 {ê}{{\^{e}}}1 {ë}{{\"e}}1
    {à}{{\`{a}}}1 {â}{{\^{a}}}1 {ä}{{\"a}}1
    {î}{{\^{i}}}1 {ï}{{\"i}}1
    {ô}{{\^{o}}}1 {ö}{{\"o}}1 {ó}{{\'{o}}}1
    {û}{{\^{u}}}1 {ù}{{\`{u}}}1 {ü}{{\"u}}1
    {ç}{{\c{c}}}1 {Ç}{{\c{C}}}1
    {É}{{\'{E}}}1 {È}{{\`{E}}}1 {Â}{{\^{A}}}1,
}

% ══════════════════════════════════════════════════════════════════════════════
\begin{document}

% ── PAGE DE GARDE ─────────────────────────────────────────────────────────────
\begin{titlepage}
  \centering
  \vspace*{1cm}

  % Bandeau coloré
  \begin{tikzpicture}
    \fill[TropicalTeal] (0,0) rectangle (16,2.5);
    \node[text=white, font=\Huge\bfseries] at (8,1.25)
      {TROPICAL MEDICAL};
  \end{tikzpicture}

  \vspace{1.2cm}

  {\LARGE\bfseries Rapport de Projet}\\[0.4cm]
  {\large Plateforme Médicale Numérique Sénégalaise}

  \vspace{1cm}
  \rule{12cm}{0.8pt}
  \vspace{0.5cm}

  \begin{tabular}{rl}
    \textbf{Auteur}       & Yakhoub Baby \\[4pt]
    \textbf{Email}        & \href{mailto:yakhoub.baby@univ-thies.sn}
                            {yakhoub.baby@univ-thies.sn} \\[4pt]
    \textbf{Établissement}& Université de Thiès \\[4pt]
    \textbf{Année}        & 2025 -- 2026 \\
  \end{tabular}

  \vspace{0.8cm}
  \rule{12cm}{0.8pt}
  \vspace{1.2cm}

  % Résumé visuel des rôles
  \begin{tikzpicture}[
    role/.style={
      draw=TropicalTeal, rounded corners=6pt,
      fill=TropicalLight, minimum width=2.2cm,
      minimum height=1cm, font=\small\bfseries,
      text=TropicalDark, drop shadow
    }
  ]
    \node[role] (p) at (0,0)     {Patient};
    \node[role] (m) at (2.6,0)   {Médecin};
    \node[role] (r) at (5.2,0)   {Réceptionniste};
    \node[role] (ph) at (8.4,0)  {Pharmacien};
    \node[role] (a) at (11.2,0)  {Admin};
  \end{tikzpicture}

  \vspace{1.5cm}

  \begin{tikzpicture}
    \fill[TropicalTeal!20] (0,0) rectangle (16,0.6);
    \node[font=\small\itshape\color{TropicalDark}] at (8,0.3)
      {Vue 3 \textbullet\ Node.js \textbullet\ Prisma 5 \textbullet\
       SQLite \textbullet\ Socket.io \textbullet\ JWT};
  \end{tikzpicture}

  \vfill
  {\small Version 4.0 — \today}
\end{titlepage}

% ── TABLE DES MATIÈRES ────────────────────────────────────────────────────────
\tableofcontents
\clearpage

% ══════════════════════════════════════════════════════════════════════════════
\chapter*{Résumé}
\addcontentsline{toc}{chapter}{Résumé}

\textbf{Tropical Medical} est une plateforme médicale numérique conçue pour
répondre aux besoins des établissements de santé au Sénégal. Elle centralise
la gestion des patients, des rendez-vous, des consultations, des dossiers
médicaux, de la pharmacie, de l'hospitalisation, de la facturation et de la
téléconsultation au sein d'une interface web moderne, accessible sur ordinateur
et téléphone mobile.

\medskip

L'application est construite autour d'une architecture découplée
\textit{frontend / backend} : le frontend est développé en \textbf{Vue 3} avec
l'API Composition, tandis que le backend repose sur \textbf{Node.js} avec le
framework \textbf{Express.js 5} et l'ORM \textbf{Prisma 5} pour la gestion
de la base de données \textbf{SQLite}. La communication en temps réel est
assurée par \textbf{Socket.io}.

\medskip

Cinq profils d'utilisateurs sont pris en charge : \textit{Patient, Médecin,
Réceptionniste, Pharmacien} et \textit{Administrateur}, chacun disposant d'un
espace personnalisé et de fonctionnalités adaptées à son rôle.

\bigskip

\begin{center}
\begin{tabular}{lc}
  \toprule
  \textbf{Indicateur} & \textbf{Valeur} \\
  \midrule
  Nombre de rôles utilisateurs   & 5 \\
  Nombre de modules fonctionnels & 14 \\
  Nombre de tables en base       & 22 \\
  Communication temps réel       & Socket.io \\
  Authentification               & JWT (JSON Web Token) \\
  \bottomrule
\end{tabular}
\end{center}

% ══════════════════════════════════════════════════════════════════════════════
\chapter{Introduction}

\section{Contexte}

Le système de santé sénégalais fait face à plusieurs défis opérationnels :
gestion manuelle des dossiers patients, manque de coordination entre les
différents acteurs médicaux (médecins, pharmaciens, réceptionnistes), et
difficulté d'accès aux soins pour les populations éloignées. La digitalisation
des processus médicaux représente un levier important pour améliorer la qualité
des soins et l'efficacité administrative.

\section{Objectifs}

L'objectif principal de ce projet est de concevoir et développer une application
web complète permettant à un établissement de santé de gérer l'ensemble de ses
activités cliniques et administratives. Plus spécifiquement, il s'agit de :

\begin{itemize}[leftmargin=2cm]
  \item Centraliser les dossiers médicaux des patients sous format numérique ;
  \item Fluidifier la prise de rendez-vous et la gestion de la file d'attente ;
  \item Faciliter les prescriptions médicales et le suivi des ordonnances ;
  \item Assurer la traçabilité de la facturation et des paiements ;
  \item Permettre la téléconsultation par vidéo ;
  \item Gérer le stock de médicaments et les lots de la pharmacie ;
  \item Offrir une messagerie interne entre les acteurs de santé ;
  \item Diffuser des actualités de santé publique.
\end{itemize}

\section{Périmètre du projet}

L'application est destinée à être déployée dans des cliniques, hôpitaux de
district et centres de santé au Sénégal. Elle prend en charge cinq catégories
d'utilisateurs, chacune avec des droits et des interfaces spécifiques.

% ══════════════════════════════════════════════════════════════════════════════
\chapter{Architecture du système}

\section{Architecture générale}

L'application suit une architecture \textbf{client-serveur découplée} en trois
couches :

\begin{enumerate}
  \item \textbf{Couche présentation} (Frontend) : application web monopage
        (\textit{Single Page Application}) développée en Vue 3.
  \item \textbf{Couche métier} (Backend) : API REST sécurisée développée en
        Node.js / Express.js.
  \item \textbf{Couche données} (Base de données) : base relationnelle SQLite
        gérée via l'ORM Prisma 5.
\end{enumerate}

\medskip

\begin{center}
\begin{tikzpicture}[
  box/.style={
    draw=TropicalTeal, rounded corners=8pt, fill=TropicalLight,
    minimum width=5cm, minimum height=1.2cm, font=\bfseries,
    text=TropicalDark, drop shadow, align=center
  },
  arr/.style={-Stealth, thick, color=TropicalTeal}
]
  \node[box] (front)  at (0, 0)   {Navigateur Web\\{\small Vue 3 — SPA}};
  \node[box] (back)   at (0,-2.5) {Serveur Node.js\\{\small Express.js 5 + Socket.io}};
  \node[box] (db)     at (0,-5)   {Base de données\\{\small SQLite via Prisma 5}};

  \draw[arr] (front.south) -- node[right, font=\small]{HTTP/REST + WS} (back.north);
  \draw[arr] (back.south)  -- node[right, font=\small]{Prisma ORM}      (db.north);
  \draw[arr] (back.north)  -- (front.south);
  \draw[arr] (db.north)    -- (back.south);
\end{tikzpicture}
\end{center}

\section{Architecture du frontend}

Le frontend est organisé selon une structure modulaire par rôle :

\begin{lstlisting}[language=bash, caption={Structure du répertoire frontend/src/}]
src/
├── assets/           # Ressources statiques
├── components/       # Composants réutilisables
│   ├── layout/       # AppSidebar, AppHeader, NotifToast
│   └── MessageriePanel.vue
├── composables/      # Hooks réutilisables (useApi, useDate)
├── layouts/          # Layouts par rôle (Patient, Medecin, ...)
├── router/           # Configuration Vue Router
├── stores/           # Stores Pinia (auth, notifications, messagerie)
├── views/            # Pages par rôle
│   ├── medecin/
│   ├── patient/
│   ├── receptionniste/
│   ├── pharmacien/
│   └── admin/
├── style.css         # Design system global
└── main.js
\end{lstlisting}

\section{Architecture du backend}

\begin{lstlisting}[language=bash, caption={Structure du répertoire backend/src/}]
src/
├── app.js                  # Point d'entrée Express + Socket.io
├── config/
│   ├── prisma.js           # Instance Prisma Client
│   └── constants.js        # Constantes métier (rôles, statuts)
├── controllers/            # Logique métier par domaine
│   ├── auth.controller.js
│   ├── patient.controller.js
│   ├── medecin.controller.js
│   ├── receptionniste.controller.js
│   ├── pharmacien.controller.js
│   ├── admin.controller.js
│   ├── notifications.controller.js
│   ├── messagerie.controller.js
│   └── actualite.controller.js
├── middlewares/
│   ├── auth.js             # Vérification JWT
│   └── roleCheck.js        # Contrôle d'accès par rôle
├── routes/                 # Définition des routes API
├── helpers/
│   └── notificationHelper.js
└── cron.js                 # Tâches planifiées (node-cron)
\end{lstlisting}

% ══════════════════════════════════════════════════════════════════════════════
\chapter{Technologies utilisées}

\section{Frontend}

\subsection{Vue.js 3}

\begin{techbox}{Vue.js 3 — Framework JavaScript}
Vue.js est un framework JavaScript progressif open-source pour la construction
d'interfaces utilisateur. La version 3, utilisée dans ce projet, introduit
l'\textbf{API Composition} (\texttt{setup}), qui permet une meilleure
organisation du code et une réutilisabilité accrue des logiques métier via
des \textit{composables}.
\end{techbox}

\medskip

Fonctionnalités Vue 3 utilisées dans le projet :
\begin{itemize}
  \item \texttt{ref()} et \texttt{computed()} pour la réactivité ;
  \item \texttt{onMounted()} / \texttt{onUnmounted()} pour le cycle de vie ;
  \item \texttt{watch()} pour réagir aux changements d'état ;
  \item \texttt{<Teleport>} pour les modales hors du DOM parent ;
  \item \texttt{<Transition>} et \texttt{<TransitionGroup>} pour les animations ;
  \item \texttt{<script setup>} — syntaxe sucre syntaxique de l'API Composition.
\end{itemize}

\subsection{Vite 8}

Vite est un outil de build ultra-rapide pour le développement web moderne.
Il exploite les modules ES natifs du navigateur en mode développement pour
éliminer le temps de bundling, ce qui se traduit par des démarrages de serveur
quasi-instantanés. En production, il utilise \textbf{Rollup} pour générer
des bundles optimisés.

\subsection{Pinia}

Pinia est le gestionnaire d'état officiel de Vue 3. Il remplace Vuex avec
une API plus simple et une meilleure intégration TypeScript. Dans ce projet,
trois stores Pinia sont utilisés :

\begin{center}
\begin{tabular}{lll}
  \toprule
  \textbf{Store} & \textbf{Fichier} & \textbf{Rôle} \\
  \midrule
  \texttt{useAuthStore}          & auth.js          & Utilisateur connecté, token JWT \\
  \texttt{useNotificationsStore} & notifications.js & Notifications temps réel \\
  \texttt{useMessagerieStore}    & messagerie.js    & Badge messages non lus \\
  \bottomrule
\end{tabular}
\end{center}

\subsection{Vue Router 4}

Vue Router assure la navigation côté client (SPA). Le routeur est configuré
avec des \textbf{routes imbriquées par rôle} et des \textbf{gardes de navigation}
(\textit{beforeEach}) qui vérifient l'authentification JWT et le rôle avant
chaque changement de page.

\subsection{Vue I18n}

Vue I18n permet l'internationalisation de l'interface. L'application supporte
le \textbf{français} et l'\textbf{anglais} avec basculement dynamique sans
rechargement de page.

\subsection{Lucide Vue Next}

Bibliothèque d'icônes SVG légère (plus de 1~000 icônes). Utilisée pour toutes
les icônes de l'interface : menus, boutons, badges, statuts.

\subsection{Socket.io Client}

Utilisé côté frontend pour recevoir les événements temps réel émis par le
serveur (nouvelles notifications, nouveaux messages de messagerie).

\subsection{@vueuse/motion}

Bibliothèque d'animations déclaratives pour Vue 3, basée sur les animations
CSS et Web Animations API. Utilisée pour les animations d'entrée des éléments.

\subsection{Tailwind CSS 4}

Framework CSS utilitaire intégré via le plugin Vite. Complète le design system
personnalisé de l'application pour certaines mises en page.

\subsection{Montserrat — Google Fonts}

Toutes les polices de l'application utilisent la famille \textbf{Montserrat}
(poids 300 à 900), importée depuis Google Fonts. Cette police sans-serif
géométrique confère à l'interface un aspect moderne et professionnel.

% ── Backend ──────────────────────────────────────────────────────────────────
\section{Backend}

\subsection{Node.js}

Node.js est un environnement d'exécution JavaScript côté serveur, basé sur
le moteur V8 de Chrome. Son modèle non-bloquant et orienté événements le rend
particulièrement adapté aux applications web nécessitant de nombreuses
connexions simultanées.

\subsection{Express.js 5}

Express.js est le framework web minimaliste le plus utilisé avec Node.js.
La version 5, adoptée dans ce projet, apporte notamment la gestion native des
erreurs asynchrones sans \textit{try/catch} explicite.

\medskip

L'API REST expose les routes suivantes :

\begin{center}
\rowcolors{2}{TropicalLight}{white}
\begin{tabular}{ll}
  \toprule
  \textbf{Préfixe} & \textbf{Domaine} \\
  \midrule
  \texttt{/api/auth}           & Authentification \\
  \texttt{/api/patient}        & Espace patient \\
  \texttt{/api/medecin}        & Espace médecin \\
  \texttt{/api/receptionniste} & Espace réceptionniste \\
  \texttt{/api/pharmacien}     & Espace pharmacien \\
  \texttt{/api/admin}          & Administration \\
  \texttt{/api/notifications}  & Notifications \\
  \texttt{/api/messagerie}     & Messagerie interne \\
  \bottomrule
\end{tabular}
\end{center}

\subsection{Prisma 5}

Prisma est un ORM (\textit{Object-Relational Mapper}) de nouvelle génération
pour Node.js. Il génère automatiquement un client TypeScript/JavaScript typé
à partir d'un schéma déclaratif (\texttt{schema.prisma}), ce qui élimine
les requêtes SQL manuelles et réduit les erreurs.

\medskip

Fonctionnalités Prisma utilisées :
\begin{itemize}
  \item Migrations automatiques (\texttt{prisma migrate dev}) ;
  \item Peuplement de la base (\texttt{prisma db seed}) ;
  \item Studio visuel (\texttt{prisma studio}) pour inspecter les données ;
  \item Relations \texttt{@relation}, clés étrangères, contraintes \texttt{@unique}.
\end{itemize}

\subsection{SQLite}

SQLite est utilisé comme base de données relationnelle. Légère et sans serveur
dédié, elle est idéale pour le développement et les petites structures.
La base de données est un seul fichier (\texttt{dev.db}) stocké localement.

\subsection{JSON Web Token (JWT)}

L'authentification repose sur les \textbf{JWT} (\textit{JSON Web Tokens}).
Lors de la connexion, le serveur génère un token signé contenant l'identifiant
et le rôle de l'utilisateur. Ce token est envoyé dans l'en-tête
\texttt{Authorization: Bearer <token>} à chaque requête et vérifié par le
middleware \texttt{auth.js}.

\subsection{Socket.io}

Socket.io implémente la communication bidirectionnelle en temps réel entre
le serveur et les clients via WebSocket (avec fallback HTTP long-polling).
Utilisations dans le projet :

\begin{itemize}
  \item Réception des \textbf{notifications} instantanées pour chaque
        utilisateur connecté (chambre \texttt{user:<id>}) ;
  \item Émission d'un événement \texttt{nouveau\_message} lors de l'envoi
        d'un message de messagerie.
\end{itemize}

\subsection{bcryptjs}

Bibliothèque de hachage des mots de passe avec l'algorithme \textbf{bcrypt}.
Tous les mots de passe sont stockés sous forme hachée (jamais en clair)
avec un facteur de coût de 10 itérations.

\subsection{node-cron}

Planificateur de tâches CRON pour Node.js. Utilisé pour exécuter des tâches
automatiques à intervalles réguliers (rappels de rendez-vous, nettoyage
de données expirées, etc.).

\subsection{dotenv}

Charge les variables d'environnement depuis le fichier \texttt{.env} dans
\texttt{process.env}. Permet de séparer la configuration sensible
(secrets JWT, URL de base de données) du code source.

% ══════════════════════════════════════════════════════════════════════════════
\chapter{Modèle de données}

\section{Vue d'ensemble}

La base de données comporte \textbf{22 tables} réparties en six domaines
fonctionnels. Le schéma est défini dans \texttt{backend/prisma/schema.prisma}.

\begin{center}
\begin{tikzpicture}[
  domain/.style={
    draw=TropicalTeal, rounded corners=6pt, fill=TropicalLight,
    minimum width=3.5cm, minimum height=1cm,
    font=\small\bfseries, text=TropicalDark, align=center, drop shadow
  }
]
  \node[domain] (u)   at (0, 4)    {Utilisateurs\\{\tiny Utilisateur, Patient, Médecin}};
  \node[domain] (rdv) at (5, 4)    {RDV \& File attente\\{\tiny RendezVous, PassageFileAttente}};
  \node[domain] (d)   at (10, 4)   {Dossier médical\\{\tiny DossierMedical, Consultation}};
  \node[domain] (ph)  at (0, 1)    {Pharmacie\\{\tiny Medicament, Ordonnance, LotMedicament}};
  \node[domain] (h)   at (5, 1)    {Hospitalisation\\{\tiny Lit, Hospitalisation}};
  \node[domain] (f)   at (10, 1)   {Facturation\\{\tiny Facture, Paiement, Transaction}};
  \node[domain] (n)   at (5,-1.5)  {Notifications \& Messagerie\\{\tiny Notification, Interaction, Actualite}};

  \draw[TropicalTeal, thick] (u) -- (rdv);
  \draw[TropicalTeal, thick] (rdv) -- (d);
  \draw[TropicalTeal, thick] (d) -- (ph);
  \draw[TropicalTeal, thick] (h) -- (f);
  \draw[TropicalTeal, thick] (u) -- (n);
\end{tikzpicture}
\end{center}

\section{Tables principales}

\begin{longtable}{p{3.5cm} p{3cm} p{7cm}}
  \toprule
  \textbf{Table} & \textbf{Domaine} & \textbf{Description} \\
  \midrule
  \endhead
  \texttt{Utilisateur}     & Auth       & Compte de base (nom, email, rôle, mot de passe haché) \\
  \texttt{Patient}         & Profil     & Données médicales du patient (groupe sanguin, allergies, antécédents) \\
  \texttt{Medecin}         & Profil     & Spécialité, numéro d'ordre, tarif consultation \\
  \texttt{RendezVous}      & Planning   & Date, type (présentiel/téléconsultation), statut, motif \\
  \texttt{PassageFileAttente} & Flux    & Gestion de la file avec niveau d'urgence \\
  \texttt{DossierMedical}  & Clinique   & Dossier central, relié aux consultations et ordonnances \\
  \texttt{Consultation}    & Clinique   & Anamnèse, examen clinique, diagnostic, traitement \\
  \texttt{Ordonnance}      & Pharmacie  & Prescriptions avec lignes de médicaments \\
  \texttt{Medicament}      & Pharmacie  & Catalogue médicaments (stock, prix, DCI) \\
  \texttt{LotMedicament}   & Pharmacie  & Traçabilité des lots (numéro, expiration, quantité) \\
  \texttt{Lit}             & Hospit.    & Numéro, chambre, étage, type, statut \\
  \texttt{Hospitalisation} & Hospit.    & Admission, sortie, coût par jour, motif \\
  \texttt{Facture}         & Finances   & Montant total, part assurance, part patient, statut \\
  \texttt{Paiement}        & Finances   & Mode de paiement, référence, montant \\
  \texttt{Teleconsultation}& Numérique  & URL de la salle, plateforme, durée \\
  \texttt{Notification}    & Alertes    & Titre, message, canal, lu/non lu \\
  \texttt{Interaction}     & Messagerie & Message, expéditeur, destinataire, lu/non lu \\
  \texttt{Actualite}       & Info santé & Titre, contenu, catégorie, auteur, publiée \\
  \bottomrule
\end{longtable}

% ══════════════════════════════════════════════════════════════════════════════
\chapter{Fonctionnalités développées}

\section{Authentification et contrôle d'accès}

Le système d'authentification repose sur :
\begin{itemize}
  \item \textbf{Inscription/Connexion} avec email et mot de passe haché par bcrypt ;
  \item Génération d'un \textbf{token JWT} valide 24~heures à la connexion ;
  \item Stockage du token dans \texttt{localStorage} côté client ;
  \item Middleware \texttt{auth.js} vérifiant la signature du token à chaque
        appel API ;
  \item Middleware \texttt{roleCheck.js} autorisant ou refusant l'accès selon
        le rôle de l'utilisateur (\texttt{PATIENT, MEDECIN, RECEPTIONNISTE,
        PHARMACIEN, ADMIN}).
\end{itemize}

\section{Tableau de bord par rôle}

Chaque rôle dispose d'un tableau de bord personnalisé avec des KPI
(\textit{Key Performance Indicators}) animés et des raccourcis vers les modules
principaux :

\begin{center}
\rowcolors{2}{TropicalLight}{white}
\begin{tabular}{ll}
  \toprule
  \textbf{Rôle} & \textbf{KPI affichés} \\
  \midrule
  Patient        & Prochains RDV, factures en attente, ordonnances actives \\
  Médecin        & File d'attente, consultations du jour, urgences \\
  Réceptionniste & RDV du jour, lits libres/occupés, paiements récents \\
  Pharmacien     & Stock faible, ordonnances à délivrer, lots expirants \\
  Admin          & Nombre d'utilisateurs par rôle, activité globale \\
  \bottomrule
\end{tabular}
\end{center}

\section{Gestion des rendez-vous}

\begin{itemize}
  \item Prise de RDV en ligne par le patient (choix médecin, date, motif) ;
  \item Types de RDV : présentiel ou téléconsultation ;
  \item Statuts : \texttt{EN\_ATTENTE}, \texttt{CONFIRME}, \texttt{ANNULE},
        \texttt{TERMINE} ;
  \item Filtres par statut avec onglets colorés.
\end{itemize}

\section{File d'attente}

Gestion en temps quasi-réel des patients en attente de consultation :
\begin{itemize}
  \item Enregistrement à l'arrivée avec niveau d'urgence (FAIBLE, MODERE,
        URGENT, CRITIQUE) ;
  \item Appel des patients par le médecin ;
  \item Conversion automatique en consultation à l'appel.
\end{itemize}

\section{Dossiers médicaux}

Chaque patient possède un dossier médical unique contenant :
\begin{itemize}
  \item Historique des consultations (anamnèse, diagnostic, traitement) ;
  \item Vaccinations avec dates de rappel ;
  \item Examens médicaux (biologie, radiologie, imagerie) ;
  \item Antécédents médicaux et chirurgicaux ;
  \item Maladies chroniques et allergies.
\end{itemize}

\section{Ordonnances et pharmacie}

\begin{itemize}
  \item Création d'ordonnances par le médecin avec sélection de médicaments
        depuis le catalogue ;
  \item Impression PDF de l'ordonnance dans une fenêtre dédiée ;
  \item Délivrance par le pharmacien avec mise à jour du stock ;
  \item Gestion des lots de médicaments (numéro, date d'expiration, quantité) ;
  \item Alertes automatiques pour les stocks faibles.
\end{itemize}

\section{Hospitalisation et gestion des lits}

\begin{itemize}
  \item Grille visuelle des lits par étage avec code couleur (Libre / Occupé /
        Réservé / Maintenance) ;
  \item Admission d'un patient sur un lit libre directement depuis la grille ;
  \item Suivi du coût par jour et calcul automatique de la facture de séjour ;
  \item Libération du lit à la sortie.
\end{itemize}

\section{Facturation et paiements}

\begin{itemize}
  \item Génération automatique d'une facture après consultation ou hospitalisation ;
  \item Calcul de la part assurance et de la part patient ;
  \item Modes de paiement : espèces, carte bancaire, Mobile Money
        (Orange Money, Wave, Free Money) ;
  \item Suivi des statuts : \texttt{EN\_ATTENTE}, \texttt{PARTIELLEMENT\_PAYE},
        \texttt{PAYE}.
\end{itemize}

\section{Téléconsultation}

\begin{itemize}
  \item Planification d'une séance de vidéo-consultation lors de la prise
        de RDV ;
  \item Génération d'un lien de salle \textbf{Jitsi Meet} ;
  \item Accès à la salle depuis le tableau de bord du médecin et du patient ;
  \item Enregistrement de la durée de la session.
\end{itemize}

\section{Messagerie interne}

Système de messagerie directe entre les acteurs de santé :
\begin{itemize}
  \item Liste des conversations avec dernière message et compteur non lus ;
  \item Envoi et réception de messages en temps (quasi) réel via
        \textbf{polling toutes les 8 secondes} ;
  \item Badge sur le menu latéral indiquant le nombre de messages non lus,
        mis à jour toutes les 15~secondes via un store Pinia ;
  \item Remise à zéro du badge à l'ouverture de la messagerie ;
  \item Accessible aux trois rôles : Patient, Médecin, Réceptionniste.
\end{itemize}

\section{Actualités santé}

\begin{itemize}
  \item Publication d'actualités de santé publique par la réceptionniste
        (campagnes de vaccination, journées mondiales, dépistages…) ;
  \item Catégories : \texttt{INFO}, \texttt{DON\_SANG}, \texttt{VACCINATION},
        \texttt{SENSIBILISATION}, \texttt{EVENEMENT}, \texttt{URGENCE} ;
  \item Consultation par le patient avec filtres par catégorie et modal
        de lecture complète ;
  \item Gestion CRUD (créer, lire, modifier, supprimer, publier/dépublier).
\end{itemize}

\section{Notifications en temps réel}

\begin{itemize}
  \item Notification in-app à chaque événement important (RDV confirmé,
        nouveau message, résultat d'examen disponible) ;
  \item Connexion Socket.io au chargement du dashboard ;
  \item Affichage sous forme de toast (\textit{pop-up} éphémère) et dans
        un panneau de notifications persistant.
\end{itemize}

% ══════════════════════════════════════════════════════════════════════════════
\chapter{Sécurité}

\section{Mesures implémentées}

\begin{center}
\rowcolors{2}{TropicalLight}{white}
\begin{tabular}{p{4cm} p{9cm}}
  \toprule
  \textbf{Mesure} & \textbf{Description} \\
  \midrule
  Hachage des mots de passe & bcryptjs avec 10 rounds de salage \\
  Authentification stateless & JWT signé avec secret configurable \\
  Contrôle d'accès & Middleware roleCheck par route API \\
  CORS configuré & Origines autorisées limitées au frontend \\
  Variables sensibles & Secrets dans \texttt{.env} (exclu du dépôt Git) \\
  Validation des entrées & Vérification côté serveur avant chaque écriture \\
  Séparation des espaces & Chaque rôle n'accède qu'à ses propres données \\
  \bottomrule
\end{tabular}
\end{center}

\section{Gestion des secrets}

Les données sensibles sont stockées dans le fichier \texttt{.env} qui n'est
\textbf{jamais versionné} (présent dans \texttt{.gitignore}). Un fichier
\texttt{.env.example} documente les variables nécessaires sans les valeurs
réelles :

\begin{lstlisting}[caption={Contenu de backend/.env.example}]
DATABASE_URL="file:./dev.db"
JWT_SECRET=remplacer_par_une_cle_secrete_longue_et_aleatoire
JWT_EXPIRES_IN=24h
PORT=5000
NODE_ENV=development
FRONTEND_URL=http://localhost:5173
\end{lstlisting}

% ══════════════════════════════════════════════════════════════════════════════
\chapter{Interface utilisateur et design}

\section{Design system}

Un design system cohérent est défini dans \texttt{frontend/src/style.css}
via des \textbf{variables CSS} (\textit{custom properties}) :

\begin{lstlisting}[language=CSS, caption={Extrait des variables CSS globales}]
:root {
  --color-primary:   #0E9F8E;   /* Teal medical  */
  --color-danger:    #EF4444;   /* Rouge urgence */
  --color-warning:   #F59E0B;   /* Ambre alerte  */
  --font-sans:    'Montserrat', system-ui, sans-serif;
  --font-display: 'Montserrat', sans-serif;
  --radius-card:  1.125rem;  /* Arrondi des cartes */
}
\end{lstlisting}

\section{Typographie}

Toute l'application utilise la police \textbf{Montserrat} (Google Fonts),
disponible en poids 300 (Light) à 900 (Black). Ce choix offre une lisibilité
optimale sur écran et une apparence moderne et professionnelle adaptée au
secteur médical.

\section{Responsive design}

L'interface s'adapte aux différentes tailles d'écran grâce aux \textit{media
queries} CSS. La barre latérale se replie automatiquement sur mobile (en
dessous de 1024~px) et un bouton hamburger permet de l'afficher en superposition.

\section{Animations}

L'application intègre plusieurs niveaux d'animations :
\begin{itemize}
  \item \textbf{Entrées en cascade} (\textit{stagger}) : les cartes apparaissent
        séquentiellement avec un délai croissant ;
  \item \textbf{Ressort} (\textit{spring}) : courbe \texttt{cubic-bezier(.34,1.56,.64,1)}
        pour un effet rebond naturel ;
  \item \textbf{Bande lumineuse} (\textit{shine}) : reflet animé sur les
        barres colorées des cartes statistiques ;
  \item \textbf{Pulse} : clignotement doux sur les badges « Occupés » pour
        attirer l'attention ;
  \item \textbf{Compteurs animés} : les chiffres des KPI s'incrémentent au
        chargement (\texttt{AnimCounter}).
\end{itemize}

\section{Mode sombre}

L'application supporte un mode sombre activable depuis les paramètres
utilisateur. Les couleurs basculent via les variables CSS en appliquant
la classe \texttt{dark} sur l'élément \texttt{<html>}.

% ══════════════════════════════════════════════════════════════════════════════
\chapter{Déploiement et gestion de version}

\section{Structure du dépôt}

\begin{lstlisting}[language=bash, caption={Structure racine du projet}]
tropical-medical/
├── backend/         # API Node.js
│   ├── src/
│   ├── prisma/
│   ├── .env.example
│   └── package.json
├── frontend/        # SPA Vue 3
│   ├── src/
│   ├── index.html
│   └── package.json
├── .gitignore       # Exclut .env, node_modules, dev.db
└── README.md
\end{lstlisting}

\section{Fichier .gitignore}

Le fichier \texttt{.gitignore} à la racine exclut du dépôt Git :
\begin{itemize}
  \item Les dossiers \texttt{node\_modules/} (dépendances) ;
  \item Les fichiers \texttt{.env} (secrets) ;
  \item La base de données SQLite (\texttt{*.db}, \texttt{*.db-journal}) ;
  \item Le client Prisma généré (\texttt{backend/generated/}) ;
  \item Le build de production (\texttt{frontend/dist/}).
\end{itemize}

\section{Commandes de lancement}

\begin{lstlisting}[language=bash, caption={Démarrage du projet en développement}]
# Backend
cd backend
npm install
npx prisma migrate dev
node prisma/seed.js    # Donnees initiales
npm run dev            # Port 5000

# Frontend (autre terminal)
cd frontend
npm install
npm run dev            # Port 5173
\end{lstlisting}

% ══════════════════════════════════════════════════════════════════════════════
\chapter{Conclusion}

\section{Bilan}

Ce projet a abouti à la réalisation d'une plateforme médicale complète et
fonctionnelle, couvrant l'ensemble du cycle de vie d'un patient au sein d'un
établissement de santé. L'application \textbf{Tropical Medical} propose :

\begin{itemize}
  \item Une architecture moderne et maintenable (séparation frontend/backend) ;
  \item Une interface intuitive adaptée à cinq profils d'utilisateurs distincts ;
  \item Des fonctionnalités temps réel via Socket.io ;
  \item Un système de sécurité robuste basé sur JWT et le contrôle d'accès
        par rôle ;
  \item Un design system cohérent avec animations et thème sombre.
\end{itemize}

\section{Perspectives d'évolution}

Plusieurs améliorations peuvent être envisagées pour les versions futures :

\begin{enumerate}
  \item Migration vers \textbf{PostgreSQL} ou \textbf{MySQL} pour la production ;
  \item Intégration d'une \textbf{API de paiement mobile} réelle (Orange Money,
        Wave) ;
  \item Module de \textbf{rapports et statistiques} avec graphiques avancés ;
  \item Application mobile native (\textbf{React Native} ou \textbf{Capacitor}) ;
  \item Déploiement sur un serveur cloud (\textbf{VPS} ou \textbf{PaaS}).
\end{enumerate}

% ══════════════════════════════════════════════════════════════════════════════
\begin{thebibliography}{99}

\bibitem{vuejs}
  Evan You et al.,
  \textit{Vue.js 3 — The Progressive JavaScript Framework},
  \url{https://vuejs.org}, 2024.

\bibitem{pinia}
  Eduardo San Martin Morote,
  \textit{Pinia — The intuitive store for Vue.js},
  \url{https://pinia.vuejs.org}, 2024.

\bibitem{prisma}
  Prisma Inc.,
  \textit{Prisma — Next-generation Node.js and TypeScript ORM},
  \url{https://www.prisma.io}, 2024.

\bibitem{express}
  OpenJS Foundation,
  \textit{Express.js — Fast, unopinionated, minimalist web framework for Node.js},
  \url{https://expressjs.com}, 2024.

\bibitem{socketio}
  Socket.io Contributors,
  \textit{Socket.IO — Bidirectional and low-latency communication},
  \url{https://socket.io}, 2024.

\bibitem{jwt}
  M. Jones, J. Bradley, N. Sakimura,
  \textit{RFC 7519 — JSON Web Token (JWT)},
  IETF, 2015.

\bibitem{vite}
  Evan You et al.,
  \textit{Vite — Next Generation Frontend Tooling},
  \url{https://vite.dev}, 2024.

\bibitem{sqlite}
  D. Richard Hipp,
  \textit{SQLite — A self-contained, serverless SQL database engine},
  \url{https://sqlite.org}, 2024.

\end{thebibliography}

\end{document}
